%%%%%%%% ICML 2026 EXAMPLE LATEX SUBMISSION FILE %%%%%%%%%%%%%%%%%

\documentclass{article}

% Font encoding for special characters
\usepackage[T1]{fontenc}

% Recommended, but optional, packages for figures and better typesetting:
\usepackage{microtype}
\usepackage{graphicx}
\usepackage{subcaption}
\usepackage{booktabs} % for professional tables
\usepackage{tabularx}
\usepackage{multirow}
\usepackage{array}

% hyperref makes hyperlinks in the resulting PDF.
% If your build breaks (sometimes temporarily if a hyperlink spans a page)
% please comment out the following usepackage line and replace
% \usepackage{icml2026} with \usepackage[nohyperref]{icml2026} above.
\usepackage{hyperref}


% Attempt to make hyperref and algorithmic work together better:
\newcommand{\theHalgorithm}{\arabic{algorithm}}

% Use the following line for the initial blind version submitted for review:
\usepackage{icml2026}

% If accepted, instead use the following line for the camera-ready submission:
% \usepackage[accepted]{icml2026}
% \usepackage[preprint]{icml2026}

% For theorems and such
\usepackage{amsmath}
\usepackage{amssymb}
\usepackage{mathtools}
\usepackage{amsthm}

% if you use cleveref..
\usepackage[capitalize,noabbrev]{cleveref}

% for colored text boxes
\usepackage{tcolorbox}

%%%%%%%%%%%%%%%%%%%%%%%%%%%%%%%%
% THEOREMS
%%%%%%%%%%%%%%%%%%%%%%%%%%%%%%%%
\theoremstyle{plain}
\newtheorem{theorem}{Theorem}[section]
\newtheorem{proposition}[theorem]{Proposition}
\newtheorem{lemma}[theorem]{Lemma}
\newtheorem{corollary}[theorem]{Corollary}
\theoremstyle{definition}
\newtheorem{definition}[theorem]{Definition}
\newtheorem{assumption}[theorem]{Assumption}
\theoremstyle{remark}
\newtheorem{remark}[theorem]{Remark}

% Todonotes is useful during development; simply uncomment the next line
%    and comment out the line below the next line to turn off comments
%\usepackage[disable,textsize=tiny]{todonotes}
\usepackage[textsize=tiny]{todonotes}


%%%%%%%%%%%%%%%%%%%%%%%%%%%%%%%%
% Custom packages
%%%%%%%%%%%%%%%%%%%%%%%%%%%%%%%%
\usepackage{algorithm}
\usepackage{algorithmic}
\usepackage{xcolor}
\usepackage{float} % for [H] float placement
\usepackage{tikz}
\usetikzlibrary{positioning, calc, arrows.meta, shapes}

% The \icmltitle you define below is probably too long as a header.
% Therefore, a short form for the running title is supplied here:
\icmltitlerunning{Global Expert Choice Networks}

\begin{document}

\twocolumn[
% \icmltitle{Global Expert Choice Models for Autoregressive Language Modeling}
\icmltitle{Global Expert Choice: Autoregressive Language Modeling with Dynamic Computation Allocation and Perfect Load Balancing}

% It is OKAY to include author information, even for blind
% submissions: the style file will automatically remove it for you
% unless you've provided the [accepted] option to the icml2026
% package.

% List of affiliations: The first argument should be a (short)
% identifier you will use later to specify author affiliations
% Academic affiliations should list Department, University, City, Region, Country
% Industry affiliations should list Company, City, Region, Country

% You can specify symbols, otherwise they are numbered in order.
% Ideally, you should not use this facility. Affiliations will be numbered
% in order of appearance and this is the preferred way.
% \icmlsetsymbol{equal}{*}

\begin{icmlauthorlist}
\icmlauthor{Ryan Sun}{lehigh}
\icmlauthor{Yixin Liu}{lehigh}
\icmlauthor{Yonghui Wu}{UF}
\icmlauthor{Lichao Sun}{lehigh}
\end{icmlauthorlist}

\icmlaffiliation{lehigh}{Computer Science and Engineering, Lehigh University, Bethlehem, PA, USA}
\icmlaffiliation{UF}{MD-HOBI-BIOMED INFORMATICS, University of Florida, Gainesville, FL, USA}
\icmlcorrespondingauthor{Lichao Sun}{lis221@lehigh.edu}

% You may provide any keywords that you
% find helpful for describing your paper; these are used to populate
% the "keywords" metadata in the PDF but will not be shown in the document
\icmlkeywords{Large Language Models, Mixture of Experts, Sparse Architectures, Expert Choice}

\vskip 0.3in
]

% this must go after the closing bracket ] following \twocolumn[ ...

% This command actually creates the footnote in the first column
% listing the affiliations and the copyright notice.
% The command takes one argument, which is text to display at the start of the footnote.
% The \icmlEqualContribution command is standard text for equal contribution.
% Remove it (just {}) if you do not need this facility.

\printAffiliationsAndNotice{}  % leave blank if no need to mention equal contribution
% \printAffiliationsAndNotice{\icmlEqualContribution} % otherwise use the standard text.

\begin{abstract}
% Mixture of Experts (MoE) models have proven effective at improving performance and efficiency of large language models (LLMs) using sparse architectures. Among their variants, Expert Choice (EC) models show great promise for dynamically allocation of computation and perfect load balancing. However, their adoption in has been limited due to lower performance and the non-causal routing. In this paper, we argue that instead of sequence-level or batch-level top-$k$ routing, EC's router decisions should be made at the distribution level, aiming for expert activations to be balanced globally rather than locally. Such a model, with a fixed cutoff threshold for each expert, immediately becomes causal as the decision for the current token is independent of future tokens. The resulting model, Global Expert Choice (GEC), achieves a better trade-off between performance and computation efficiency.
% Unlike Token Choice MoE (TC-MoE) models, where each token selects its top experts, Expert Choice (EC) models have experts select the top tokens within each batch, achieving perfect load balancing and avoiding expert collapse. However, EC adoption in LLMs has been limited due to non-causal routing: determining which tokens an expert processes requires information about future tokens. To address this, we propose Global Expert Choice (GEC), which preserves EC's load-balancing benefits while restoring causality. Rather than selecting top tokens within a batch, GEC aims to pick them from the whole population by maintaining an exponential moving average (EMA) of each expert's selection threshold, estimated from historical batches. At both training and inference time, each token is independently routed based on whether its routing score exceeds this threshold—eliminating dependence on future tokens and minimizing train-inference mismatch. Extensive experiments show that GEC greatly outperforms TC-MoE models under same settings. 

% Mixture of Experts (TC-MoE) routes each token to a fixed number of experts, limiting dynamic computation allocation and suffering from load imbalance. Expert Choice (EC) inverts this by having each expert select its top tokens within a batch but breaks of causality, making it unusable for autoregressive language modelling. To tackle these problems, we propose Global Expert Choice models. We extend the routing pool from local batch to the whole population by maintaining an exponential moving average (EMA) of each expert's selection threshold. At both training and inference time, each token is independently routed based on whether its routing score exceeds this threshold—eliminating dependence on future tokens and minimizing train-inference mismatch. On small scale LLM pretraining tasks, we show that GEC significantly outperforms TC-MoE by 0.05 in cross entropy loss. 

Token-choice Mixture of Experts (TC-MoE) routes each token to a fixed number of experts, limiting dynamic computation allocation and suffering from load imbalance. Expert Choice (EC) solves the issues by having each expert select its top tokens within a batch. However, EC breaks causality, making it unsuitable for autoregressive language modeling. To tackle these problems, we propose Global Expert Choice (GEC). Compared to EC, we extend the routing pool from the local batch to the global token distribution by maintaining an exponential moving average (EMA) of each expert's selection threshold. At both training and inference time, each token is independently routed based on whether its routing score exceeds this threshold—eliminating dependence on future tokens and minimizing train-inference mismatch.
In pretraining experiments scaling to 2.4B parameters on FineWeb-Edu, GEC achieves 0.067 lower cross-entropy loss than TC-MoE.
\end{abstract}

\section{Introduction}

\begin{figure*}[t]
\centering
% Placeholder illustration commented out for archival
% \fbox{\parbox{0.95\textwidth}{\centering\vspace{3cm}\textbf{Placeholder: GEC Illustration}\\\vspace{0.5cm}Left: EC selects top-$k$ tokens within batch (requires future tokens)\\Right: GEC uses EMA threshold from historical batches (causal)\vspace{3cm}}}
\includegraphics[width=0.95\textwidth]{figures/illustration/gec.pdf}
\caption{Comparison of TC, EC and GEC routing. \textbf{Left:} TC-MoE selects the top-$G$ experts for each token, causing load imbalance. \textbf{Middle:} EC selects the top-$k$ tokens for each expert within a batch, requiring access to all tokens including future ones. \textbf{Right:} GEC maintains an EMA of the top-$k$ cutoff threshold from historical batches and routes tokens via binary thresholding, enabling fully causal routing.}
\label{fig:gec_illustration}
\end{figure*}

% \input{tikz/illustration.tex}

Token Choice Mixture of Experts architectures (TC-MoE)~\cite{shazeer2017outrageously, lepikhin2021gshard, fedus2022switch} have emerged as a leading approach to scale language models efficiently, powering frontier models like DeepSeek-V3~\cite{deepseekv3}. Instead of activating all parameters in the Multi Layer Perceptron (MLP), TC-MoE routes each token to a fixed number of experts (top-$G$) by routing score, enabling expert specialization and increased model capacity while keeping compute nearly constant. However, this fixed-$G$ constraint causes two key limitations. First, routing decisions are made independently per token with a fixed compute budget, leaving no room for dynamically allocating more computation to harder tokens. Second, load balancing remains difficult: popular experts become overloaded while others are underutilized, leading to expert collapse~\cite{chi2022representation} and inefficient hardware utilization. Existing approaches address load imbalance through auxiliary losses~\cite{fedus2022switch, zoph2022stmoe} or PID controllers~\cite{longcat2025, wang2024auxiliarylossfreeloadbalancingstrategy}, but these add training complexity and often struggle to achieve true balance.

In comparison, sequence-level Expert Choice (Seq-EC) models~\cite{zhou2022mixture} invert the routing direction: instead of tokens selecting experts, each expert selects its top-$k$ tokens by routing score within a sequence. Since every expert processes exactly $k$ tokens, Seq-EC achieves perfect load balancing by construction, eliminating the need for auxiliary losses or controllers. Moreover, because a token may be selected by multiple experts (or none), Seq-EC naturally enables dynamic computation allocation—harder tokens can receive more expert attention. Seq-EC has been successfully applied to bidirectional models including diffusion~\cite{sun2024ecdit, shi2025diffmoe} and multimodal architectures~\cite{lin2024moma, zong2024mova} with observed higher training throughput~\cite{muennighoff2024olmoe}. However, Seq-EC violates causality: selecting a token requires comparing it against all other tokens in the sequence—including future tokens unavailable during autoregressive generation. 
Moreover, the model can cheat during training by exploiting information leaked from future tokens through routing decisions~\cite{wang2024auxiliarylossfreeloadbalancingstrategy}.
% This creates two problems: (1) during training, the model can ``cheat'' by peeking at future tokens for routing decisions~\cite{wang2024auxiliarylossfreeloadbalancingstrategy}, and (2) at inference, the required future context does not exist. 

Recent work addresses this by extending EC to the batch level, performing top-$k$ across all tokens in a batch rather than per sequence.
Empirically, larger batch sizes improve EC performance~\cite{ni2025openmoe2, shi2025diffmoe} while reducing future information leakage~\cite{wang2024auxiliarylossfreeloadbalancingstrategy}.
Batch-level EC also enables (though still awkwardly) autoregressive generation by selecting top-$k$ across current tokens from different user requests~\cite{ludziejewski2024scaling}.
However, causality is not fully restored: routing still depends on other tokens in the batch, and train-inference mismatch persists when batch size and composition differ between training and deployment. 

% In this paper, we further extend the idea and push load balancing to the population level. We propose Global Expert Choice (GEC), a novel approach that addresses the causality issue while preserving EC's load-balancing benefits. Instead of selecting top-$k$ tokens within each batch, GEC routes tokens based on whether they fall in the top percentile of the router score distribution across the training population. To estimate this percentile threshold, we maintain an exponential moving average (EMA) of the batch-level top-$k$ cutoff. At both training and inference, a token independently activates an expert if its routing score exceeds this EMA threshold. Since this threshold is determined solely by past data, routing becomes fully causal, with no train-inference mismatch. Moreover, our EMA mechanism can also enable causal inference for existing EC models without retraining.
In this paper, we push load balancing to the population level and propose Global Expert Choice models (GEC). Instead of selecting top-$k$ tokens within each batch, GEC maintains an exponential moving average (EMA) of the batch-level top-$k$ cutoff as an estimate of the population-level threshold. At both training and inference, a token is independently activated if its routing score exceeds this threshold. Since the threshold is determined solely by past data, routing becomes fully causal with no train-inference mismatch. Moreover, our EMA mechanism can enable causal inference for existing EC models without retraining. 

% In comparison, Expert Choice (EC)~\cite{zhou2022mixture} inverts the routing direction: instead of tokens selecting experts, each expert selects its top-$k$ tokens by routing score within a sequence or batch. Since every expert processes exactly $k$ tokens, EC achieves perfect load balancing by construction, eliminating the need for auxiliary losses or controllers. Moreover, because a token may be selected by multiple experts (or none), EC naturally enables dynamic computation allocation—harder tokens can receive more expert attention. EC has been successfully applied to bidirectional models including diffusion~\cite{sun2024ecdit, shi2025diffmoe} and multimodal architectures~\cite{lin2024moma, zong2024mova}.

% However, EC violates causality: selecting a token requires comparing it against all other tokens in the batch—including future tokens unavailable during autoregressive generation. This creates two problems: (1) during training, the model can ``cheat'' by peeking at future tokens for routing decisions~\cite{wang2024auxiliarylossfreeloadbalancingstrategy}, and (2) at inference, the required future context does not exist. Prior workarounds include training auxiliary predictors~\cite{raposo2024mixture} or performing top-$k$ across batched user requests~\cite{ludziejewski2024scaling}, but these add complexity and leave training-time issues unresolved (see Section~\ref{subsec:causal_generation_ec}). 

% In this paper, we propose Global Expert Choice (GEC), a novel approach that addresses the causality issue while preserving EC's load-balancing benefits. Instead of selecting top-$k$ tokens within each batch, GEC routes tokens based on whether they fall in the top percentile of the router score distribution across the training population. To estimate this percentile threshold, we maintain an exponential moving average (EMA) of the batch-level top-$k$ cutoff. At both training and inference, a token independently activates an expert if its routing score exceeds this EMA threshold. Since this threshold is determined solely by past data, routing becomes fully causal, with no train-inference mismatch. Moreover, our EMA mechanism can also enable causal inference for existing EC models without retraining.

% Our contributions are twofold:
% \begin{itemize}
%     \item We solve the fundamental causality issue in Expert Choice models by proposing a global routing mechanism that maintains causality while achieving expert-level load balancing.
%     \item We demonstrate through extensive experiments that GEC achieves competitive performance with EC while enabling causal generation.
% \end{itemize}

Pretraining a 2.4B (0.56B active) language model, GEC outperforms TC by 0.067 in cross-entropy loss while achieving near-perfect load balancing.
We further show that EC's performance improves with batch size, and that models trained with large-batch EC can perform causal inference using our threshold-based routing without retraining.
%We also show that a model trained with EC with large enough batch size can perform inference using our threshold-based routing mechanism without retraining.

\section{Related Work}

\subsection{Mixture of Experts}

Mixture of Experts (MoE) scales model capacity by routing each token to a small subset of experts while keeping compute nearly constant. A learned gate selects top-$G$ experts per token~\cite{shazeer2017outrageously}, with auxiliary losses to balance load across experts~\cite{lepikhin2021gshard}. The Switch Transformer~\cite{fedus2022switch} sets $G{=}1$ for efficiency. Recent LLMs further adopt fine-grained MoE with many small experts and shared experts that remain always active to capture global knowledge~\cite{dai2024deepseek}. We incorporate shared experts in our design.

\subsection{Load Balancing}

A critical challenge in MoE systems is load balancing: without explicit constraints, routers often favor a small subset of experts, leading to expert collapse. The standard approach uses an auxiliary loss $\mathcal{L}_{\text{aux}}$ to encourage uniform expert assignment~\cite{lepikhin2021gshard, fedus2022switch}. For a batch of $N$ tokens and $GE$ experts, we define the normalized load $f_i$ and average routing probability $P_i$ for expert $i$:
\begin{align}
f_i &= \frac{E}{N} \sum_{t=1}^N z_{t,i}, \\
P_i &= \frac{1}{N} \sum_{t=1}^N p_{t,i}
\end{align}
where $z_{t,i} \in \{0,1\}$ is the indicator that expert $i$ is selected for token $t$, $p_{t,i}$ is the routing probability, and $E$ is the expansion factor (total parameters / active parameters). Under Top-$G$ routing (where granularity $G$ denotes the number of experts selected per token, i.e., $\sum_i z_{t,i} = G$), $f_i=1$ implies perfect load balancing. The auxiliary loss is then defined as:
\begin{equation}
\mathcal{L}_{\text{aux}} = \alpha \sum_{i=1}^{GE} f_i P_i
\end{equation}
where $\alpha$ is a coefficient (typically $10^{-2}$ to $10^{-4}$).
However, in distributed training, the small local batch size $N$ causes high variance in load estimation.
Global-batch load balancing~\cite{qiu2025demonsdetailimplementingload,qwen3} addresses this by computing balance statistics across all devices, providing more accurate load estimates and more meaningful gradients, leading to improved expert specialization.

Despite auxiliary losses, load imbalance persists. GShard's \emph{expert capacity}~\cite{lepikhin2021gshard} caps tokens per expert (typically 1.25--2.0$\times$ average), preventing saturation but wasting compute when underutilized and dropping tokens when overloaded. MegaBlocks~\cite{gale2023megablocks} eliminates this trade-off with dynamic sparse kernels.

Recent work explores auxiliary-loss-free alternatives. DeepSeekMoE~\cite{dai2024deepseek} introduces expert-specific bias terms $b_i$ that dynamically adjust based on load statistics. Expert selection uses biased scores $r_{t,i} + b_i$, while gating weights use original scores $r_{t,i}$, preserving specialization. The bias updates follow $b_i \leftarrow b_i + u \cdot \text{sign}(1 - f_i)$, where $f_i$ is a normalized load statistic for expert $i$ (equal to 1 under perfect balance). This eliminates the trade-off between load balancing and task performance inherent in auxiliary loss methods. LongCat-Flash~\cite{longcat2025} adopts a similar framework but replaces the sign-based update with proportional control: $\Delta b_i = u \cdot (1 - f_i)$. While DeepSeek's approach applies constant-magnitude corrections regardless of imbalance severity, proportional updates scale with the load deviation, enabling smoother convergence. 

Global Expert Choice (GEC) combines the above ideas. Instead of a per-batch top-$k$ selection for the original EC, we extend Qwen's philosophy to compute balance statistics across the entire pretrain population by maintaining a global cutoff threshold using EMA. Such number, surprisingly, functions similarly to the bias term for loss-free load balancing. We further prevent GPU overload using expert capacity. See Table~\ref{tab:conceptual_mapping} for more details.

\begin{table}[t]
\centering
\caption{Taxonomy of load balancing methods by scope (Aux loss~\cite{lepikhin2021gshard}; Global LBL~\cite{qiu2025demonsdetailimplementingload}; LossFree~\cite{wang2024auxiliarylossfreeloadbalancingstrategy}; Seq EC~\cite{zhou2022mixture}; Batch EC~\cite{ludziejewski2024scaling}).}
\label{tab:load_balance_taxonomy}
\small
\begin{tabularx}{\columnwidth}{@{}>{\centering\arraybackslash}X>{\centering\arraybackslash}X>{\centering\arraybackslash}X@{}}
\toprule
\textbf{Micro Batch/Seq} & \textbf{Batch} & \textbf{Population} \\
\midrule
Aux loss & Global LBL & -- \\
-- & LossFree & -- \\
Seq EC & Batch EC & GEC (ours) \\
\bottomrule
\end{tabularx}
\end{table}

\begin{table}[t]
\centering
\caption{Conceptual connections between GEC and recent work (LossFree~\cite{wang2024auxiliarylossfreeloadbalancingstrategy}; GShard~\cite{fedus2022switch}).}
\label{tab:conceptual_mapping}
\small
\begin{tabularx}{\columnwidth}{@{}>{\raggedright\arraybackslash}X>{\raggedright\arraybackslash}X>{\raggedright\arraybackslash}X@{}}
\toprule
\textbf{GEC} & \textbf{Similar To} & \textbf{Connection} \\
\midrule
Cutoff-EMA $c_i$ & LossFree bias $b_i$ & Per-expert scalar; no aux loss \\
\midrule
$1 - \beta$ & LossFree $\mu$ & Update rate \\
\midrule
Capacity & GShard capacity & Train/inf mismatch \\
\bottomrule
\end{tabularx}
\end{table}

\subsection{Dynamic Computation}

Dynamic computation methods aim to adaptively allocate computational resources based on input complexity. Expert Choice (EC) routing~\cite{zhou2022mixture} reverses traditional token-choice MoE: instead of each token selecting top-$k$ experts, each expert selects the top-$k$ tokens with highest routing scores (this token-choice vs.\ expert-choice distinction is also used in recent router taxonomies and empirical comparisons in vision MoEs~\cite{liu2024routersvision}). Given router logits $r \in \mathbb{R}^{N \times GE}$ for $N$ tokens (sequence/batch) and $GE$ experts, expert $i$ sets the binary dispatch indicator $z_{t,i}$ as
$$z_{t,i} = \mathbf{1}\{t \in \text{TopK}(\{r_{1,i}, \ldots, r_{N,i}\}, k)\}$$
where $k = N/E$ is the number of activated tokens for the expert.
This ensures perfect load balancing (each expert processes exactly $k$ tokens) while allowing variable computation per token (0 to $GE$ experts). Expert choice has been successfully applied to upcycling dense checkpoints into sparse MoE models~\cite{komatsuzaki2023sparse} or skipping the attention layer too~\cite{raposo2024mixture}. It also sees adoption in vision models~\cite{liu2024routersvision}, diffusion models~\cite{sun2024ecdit, shi2025diffmoe}, multimodal models (e.g., MoMa~\cite{lin2024moma}), and language diffusion models OpenMoE2~\cite{ni2025openmoe2}. Related variants further expand the expert-choice design space beyond the original formulation (e.g., ternary expert choice in TC-MoE~\cite{yan2025tcmoe}).
However, EC has a causality problem for autoregressive generation: determining $z_{t,i}=1$ requires future routing scores $r_{(t+1),i}, \ldots, r_{N,i}$, which are unavailable during inference. See Section~\ref{subsec:causal_generation_ec} for more details.

Besides EC, other approaches to dynamic computation rely on other explicit designs. ReMoE~\cite{wang2025remoe} replaces discrete TopG routing with fully differentiable ReLU-based routing and adaptive L1 regularization.
Other works~\cite{jin2024moe++, longcat2025, zeng2024adamoe} introduce zero-computation experts (e.g., zero, copy, and constant) that allow tokens to skip expert computation entirely, an approach \citet{kilian2026datasparsity} extend to multimodal modeling.
Top-P routing~\cite{liu2024unimoeaudio, jin2025dtop_p, huang2024hardertasksneedexperts, wang2025hmoe} selects experts based on cumulative probability mass, adapting expert count to routing confidence—high-confidence tokens use fewer experts while uncertain ones activate more.
Closely related confidence-adaptive compute knobs also include threshold-based routing (e.g., XMoE~\cite{yang2024xmoe}) and auto-tuning methods like DynMoE~\cite{guo2025dynamicmixtureexpertsautotuning} that let each token determine how many experts to activate while reducing sensitivity to MoE hyperparameters. Beyond MoE routing itself, conditional computation can also be applied to other Transformer components and long-context settings, e.g., CoLT5~\cite{ainslie-etal-2023-colt5}.
Early exit methods~\cite{xin2020deebert} enable sample-level dynamics by allowing tokens to exit at intermediate layers.

\subsection{Causal Generation of Expert Choice Models}
\label{subsec:causal_generation_ec}

EC poses a causality challenge: token selection requires ranking against future tokens, which are unavailable in autoregressive generation.
Prior work addresses this issue in three main ways. Predictor-based methods train an auxiliary predictor or learn per-expert thresholds to approximate oracle top-$k$ decisions, enabling causal routing at inference~\cite{raposo2024mixture, shi2025diffmoe}. Alternatively, top-$k$ selection across the current tokens from different sequences preserves causality within each sequence~\cite{ludziejewski2024scaling, wen2025seqtopk}. Recent work changes routing granularity: Lory routes at the segment level, using the previous segment to determine the next~\cite{zhong2024lory}, while SeqTopK shifts expert budgets to sequence-level selection with an \emph{Expert Cache} for autoregressive decoding~\cite{wen2025seqtopk}.
All above approaches have significant drawbacks: predictions can be noisy and unstable, and batch-level top-$k$ can impose inference-time topology constraints, leading to a large train--inference mismatch; moreover, routing that depends on global batch composition can be sensitive to batch size/composition and raises privacy/safety concerns in multi-tenant settings~\cite{wen2025seqtopk}. In contrast to EC, GEC reduces to a simple threshold test (whether token logit $r_{t,i}$ is higher than cutoff EMA $c_i$) at inference time, thus eliminating the train--inference discrepancy.

\subsection{From Batch to Population Level Statistics}

The progression from sample, batch, to population-level statistics is a recurring theme in deep learning. While techniques like Batch Normalization~\cite{ioffe2015batch} and contrastive learning~\cite{radford2021learning} rely on batch statistics, momentum-based approaches~\cite{he2020momentum,caron2021emerging} and adaptive optimizers like Adam~\cite{kingma2015adam} use Exponential Moving Averages (EMA) to approximate population distributions. 
GEC applies this principle to routing via EMA-based cutoffs.

\section{Methods}

% In this section, we introduce the Global Expert Choice (GEC) model, a novel approach that addresses the causality issue while preserving the benefits of expert choice routing.

Suppose we have an expansion rate of $E$ (i.e., the sparsity ratio). For each expert, GEC seeks to pick the top $1/E$ fraction of tokens from the full router logit distribution, rather than from a single batch. To do so, let $N$ be the number of tokens in a batch and $k = N/E$ be the ideal number of selects tokens in a batch under perfect load balance for each expert. GEC records the exponential moving average (EMA) of the cutoff threshold, i.e.\ the value of the $k$-th largest router logit of each batch.
Conceptually, this cutoff EMA $c_i$ serves as a statistical estimator of the $1/E$-quantile of the router logit distribution.
Then, for both training and inference, we route tokens via binary thresholding, setting $z_{t,i}=\mathbf{1}\{r_{t,i}>c_i\}$ where $z_{t,i}\in \{0,1\}$ is the binary indicator of whether token $t$ is routed to expert $i$.

Population-level thresholding decouples routing decisions from batch composition. Essentially, we no longer fix the number of activated tokens in a batch, but only seek to do so asymptotically. In exchange, we smooth the batch-to-batch fluctuations of the cutoff threshold, yielding more consistent routing during training. The larger effective decision pool creates greater potential for expert specialization~\cite{qiu2025demonsdetailimplementingload}. Meanwhile, routing at inference becomes fully causal as we only need the global statistics of the cutoff EMA and not any information about the current batch. Thus, we eliminate the train-inference mismatch of EC.

\begin{algorithm}[t]
\caption{Global Expert Choice Routing}
\label{alg:gec}
\begin{algorithmic}[1]
\STATE \textbf{Input:} router logits $r \in \mathbb{R}^{N \times GE}$, cutoff-EMA $\{c_i\}$, decay rate $\beta$, target selection size $k= N/E$
\FOR{expert $i = 1, \ldots, GE$}
    \STATE $z_{t,i} \gets \mathbf{1}\{r_{t,i} > c_i\} \quad \forall t$ 
    \IF{\textsc{Training}}
        \STATE $c_i \gets \beta c_i + (1-\beta)\cdot \text{kth-largest}(\{r_{t,i}\}_{t=1}^{N}, k)$
    \ENDIF
\ENDFOR
\STATE \textbf{Return} $z$, $\{c_i\}$
\end{algorithmic}
\end{algorithm}

\paragraph{Warmup.} At the beginning of training, the router logits' distribution is not stable yet. The cutoff-EMA requires several thousand steps to converge to a meaningful estimate of the population quantile. During this period, incorrect thresholds cause severe expert starvation—most tokens fail to exceed the threshold, leaving experts underutilized. To address this cold-start problem, we use standard EC routing for the first 4k steps before switching to GEC. This allows the cutoff-EMA to accumulate stable statistics under controlled load balance.
% As shown in Figure~\ref{fig:warmup}, the warmup phase eliminates the early starvation spike and yields smoother expert usage throughout training.

\paragraph{Expert Capacity.} With GEC, the number of tokens activating each expert is no longer fixed but fluctuates slightly around the target capacity. To avoid GPU out-of-memory, we enforce expert capacity during training time to bound the number of tokens activating each expert~\cite{fedus2022switch}. This design introduces a small discrepancy between training and inference time, which we find acceptable.

\paragraph{Shared Expert.} Unlike Token-choice MoE which always activates some experts, Expert Choice models may have zero experts selected for certain tokens, which we hypothesize to cause issues with later layers' routing. To address this, we add a shared expert $\mathcal{M}_0$~\cite{dai2024deepseek, komatsuzaki2023sparse} that processes every token, ensuring non-trivial expert output even when no routed experts are activated.
% \footnote{Alternatively, one may use Sinkhorn routing~\cite{zhou2022mixture} to ensure non-zero expert activation. We leave this for future work.}

\paragraph{Gating.} Following LossFree~\cite{wang2024auxiliarylossfreeloadbalancingstrategy} and Mixture-of-Depths~\cite{raposo2024mixture}, we use sigmoid gates ($p_{t,i} = \sigma(r_{t,i})$) instead of softmax gates.

\section{Experiments}
\label{sec:experiments}

\subsection{Experiment Setup}
We evaluate our methods on Nanochat~\citep{karpathy2025nanochat}, an open-source codebase for training GPT-2-like models. 
We conduct experiments at two scales: a d12 model (575M parameters, 195M active) with 12 transformer layers, and a d20 model (2.4B parameters, 561M active) with 20 transformer layers.
For MoE layers, we use 16 routed experts ($G{=}2$, $E{=}8$) plus 1 shared expert; each token activates the shared expert and on average 1 routed expert.
The first layer is kept dense following common practice~\citep{deepseekv3, wang2024auxiliarylossfreeloadbalancingstrategy} to allow meaningful routing. We train on 10B and 11.2B tokens for d12 and d20, respectively, from the FineWeb-Edu 100B dataset~\citep{penedo2024the} with a batch size of 0.5M tokens (for d20, we halve the minibatch size and use 2-step gradient accumulation). We report CE loss and CORE benchmark results~\citep{li2024dclm}. Architecture, training, and evaluation details are in Appendices~\ref{app:arch_diff}, \ref{app:training_setup}, and~\ref{app:coreeval}.


\subsection{Main Results}
We compare Global Expert Choice (GEC) against Expert Choice (EC) and Token Choice (TC) routing. All variants share the same architecture and parameter count. For GEC, we use EMA decay $\beta = 0.999$ and capacity $C = 0.5$; we also evaluate a variant with EC warmup for the first 4k steps. For EC, we sweep the global selection batch size from 2k to 512k tokens. For TC, we report variants with no load balancing, auxiliary loss ($\alpha{=}0.001$), and loss-free load balancing ($u{=}0.005$). Tables~\ref{tab:main_results} and~\ref{tab:d20_results} summarize results.
GEC+warmup consistently outperforms TC in both CE loss (by 0.05 on d12 and 0.067 on d20) and CORE (by 1.89 on d12 and 2.83 on d20). 
EC with large batch sizes achieves comparable CE loss to GEC, confirming that explicit large-batch selection and EMA-based thresholding reach similar training loss. EC 512k slightly edges out GEC+warmup on CORE (19.94 vs.\ 19.88) in d12, though both substantially outperform TC. Critically, GEC without warmup suffers degraded downstream performance (CORE 16.87, below even baseline TC at 17.98) despite matching GEC+warmup in CE loss. 



\begin{table}[t]
\centering
\caption{Main results comparing Expert Choice (EC), Token Choice (TC), and Global Expert Choice (GEC) routing. \textbf{Batch}: token pool size for routing—EC uses global selection batch, TC uses per-step batch, GEC reports effective EMA pool size $N/(1-\beta)$. TC variants: no load-balancing, auxiliary loss ($\alpha{=}0.001$), or loss-free ($u{=}0.005$). GEC uses capacity $C{=}0.5$; ``warmup'' applies EC for first 4k steps. We report validation cross-entropy (CE) loss ($\downarrow$) and CORE Eval score ($\uparrow$).}
\label{tab:main_results}
\small
\setlength{\tabcolsep}{4pt}
\begin{tabularx}{\columnwidth}{@{}>{\raggedright\arraybackslash}X c r r@{}}
\toprule
\textbf{Method} & \textbf{Batch} & \textbf{CE loss ($\downarrow$)} & \textbf{CORE ($\uparrow$)} \\
\midrule
TC  & --- &  2.893 & 17.983 \\
TC aux & 64k & 2.892 & 15.894 \\
TC loss-free & 512k & 2.898 & 18.031 \\
\midrule
EC & 2k & 2.910 & 17.91 \\
EC & 8k & 2.845 & 18.83 \\
EC & 64k & \textbf{2.841} & 18.754 \\
EC & 512k & \textit{2.843} & \textbf{19.94} \\
\midrule
GEC ($\beta{=}0.999$) & $\sim$500M & 2.844 & 16.867 \\
% GEC ($\beta{=}0.997$) & $\sim$165M & 2.845 & 16.511 \\
GEC ($\beta{=}0.999$+warmup) & 0.5M $\rightarrow$ 500M & 2.844 & \textit{19.876} \\
\bottomrule
\end{tabularx}
\end{table}

% d20 table: 
\begin{table}[t]
\centering
\caption{d20 results. We report validation cross-entropy (CE) loss ($\downarrow$) and CORE Eval score ($\uparrow$).}
\label{tab:d20_results}
\small
\setlength{\tabcolsep}{4pt}
\begin{tabularx}{\columnwidth}{@{}>{\raggedright\arraybackslash}X c r r@{}}
\toprule
\textbf{Method} & \textbf{Batch} & \textbf{CE loss ($\downarrow$)} & \textbf{CORE ($\uparrow$)} \\
\midrule
dense & --- & 2.751 & 20.43 \\
TC aux & 32k & 2.687 & 22.31 \\
EC & 256k & \textit{2.621} & \textit{24.98} \\
GEC ($\beta{=}0.999$+warmup) & 256k $\rightarrow$ 500M & \textbf{2.620} & \textbf{25.14} \\
\bottomrule
\end{tabularx}
\end{table}

\subsection{Analysis}
\label{sec:analysis}

We analyze key aspects of Global Expert Choice routing: batch size scaling, train-evaluation gap, cutoff-usage tradeoff, expert specialization, capacity constraints, and comparison to Token Choice.

\subsubsection{Batch Size Scaling}
\label{sec:batch_size_scaling}

We examine how EC's routing batch size affects model quality. Table~\ref{tab:main_results} shows EC across four batch sizes (2k, 8k, 64k, 512k tokens). Larger batches yield better performance: training CE loss improves from 2.874 (2k) to 2.855 (8k) to 2.836 (64k), with CORE Eval scores following a similar trend (17.91 → 18.76 → 19.62). This is expected: top-$k$ selection over larger token pools better approximates the population-level routing decision. However, performance saturates around 64k tokens—increasing to 512k provides no further gain (2.843 CE, 19.94 CORE Eval).

Figure~\ref{fig:ec_bsz_scaling} visualizes this scaling behavior. Notably, GEC with warmup achieves comparable performance (2.844 CE, 19.52 CORE Eval) without requiring batch size coordination, making it practical for autoregressive inference where only single tokens are available.

\begin{figure}[t]
    \centering
    \includegraphics[width=\linewidth]{figures/figures/analysis/ec_vs_gec/batch_size_sensitivity/loss_vs_bsz.pdf}
    \caption{\color[HTML]{333333}EC performance across routing batch sizes. \textcolor[HTML]{3358C4}{Training CE loss} (left axis) decreases and \textcolor[HTML]{E07800}{CORE Eval score} (right axis) increases with larger batches, but gains diminish beyond 64k tokens. GEC achieves comparable performance without batch size constraints.}
    \label{fig:ec_bsz_scaling}
\end{figure}

\begin{figure}[h]
    \centering
    \includegraphics[width=\linewidth]{figures/figures/analysis/ec_vs_gec/train_eval_discrepancy/train_vs_eval_v2.pdf}
    \caption{Train-eval gap (eval loss $-$ train loss EMA) for EC at different batch sizes. EC (2k) shows a large gap that spikes during training; EC (512k) remains near zero, demonstrating EC's sensitivity to routing batch size.}
    \label{fig:train_vs_eval}
\end{figure}

\subsubsection{Train-Evaluation Gap}
\label{sec:train_inference}

A key concern for Expert Choice is the discrepancy between training and inference. During training, EC selects the top-$k$ tokens for each expert within a batch; during autoregressive inference, future tokens are unavailable, making batch-level top-$k$ selection impossible. Prior work addresses this via auxiliary predictors~\cite{raposo2024mixture} or batch-level approximations~\cite{ludziejewski2024scaling}, but these introduce additional complexity or latency.

Our results demonstrate that this concern depends critically on the routing batch size. As shown in Table~\ref{tab:main_results}, EC with large batch sizes (64k, 512k) achieves validation loss nearly identical to GEC (2.841--2.843 vs 2.844), with comparable CORE Eval scores. However, \textbf{smaller batch sizes reveal significant train-inference mismatch}: EC at 2k tokens shows degraded CORE Eval performance (17.91 vs 19.94 at 512k) and evaluation loss (2.910 vs 2.843). This gap arises because top-$k$ selection over a small batch is a noisy estimate of the population-level routing decision; at inference (batch size 1), this noise becomes extreme.

Figure~\ref{fig:train_vs_eval} illustrates this: EC (2k) shows a large train-evaluation gap, while EC (512k) show close alignment between train loss EMA and eval loss. GEC's cutoff-EMA mechanism addresses this by maintaining a population-level threshold that is independent of batch size, enabling consistent routing at inference without large-batch coordination.


% The warmup strategy---using EC's top-$k$ routing for the first 4k steps before switching to GEC's threshold routing---further demonstrates compatibility. The model smoothly transitions from top-$k$ to threshold-based selection without performance degradation, confirming that the cutoff-EMA accurately captures the routing boundary learned during top-$k$ training.


\subsubsection{Cutoff vs Expert Usage Tradeoff}
\label{sec:cutoff_usage}

EC and GEC achieve routing stability through complementary mechanisms. EC enforces \textbf{fixed expert usage}: each expert selects exactly top-$k$ tokens, guaranteeing usage of $1/E$ per expert. However, the cutoff threshold varies batch-to-batch, with standard deviation scaling as $\mathcal{O}(1/\sqrt{N})$. 

GEC inverts this: the cutoff-EMA provides a \textbf{stable threshold} ($\beta = 0.999$), while expert usage fluctuates around the capacity target. Figure~\ref{fig:cutoff_vs_usage} shows EC has a non-zero absolute deviation in per batch cutoff and the cutoff-EMA, while GEC achieves zero deviation by design. This enables consistent inference without large-batch coordination.

\begin{figure}[t]
    \centering
    \includegraphics[width=\linewidth]{figures/figures/analysis/ec_vs_gec/ema_cutoff_usage_tradeoff/cutoff_vs_usage.pdf}
    \caption{Cutoff stability vs expert usage tradeoff. \textbf{Top:} EC's cutoff deviates from its EMA even with 512k batch size. GEC's cutoff equals its EMA by design (zero deviation). \textbf{Bottom:} EC usage is fixed at $1/16$ by top-$k$ selection; GEC usage varies around the capacity target. This complementary tradeoff enables GEC's causal inference without large batches.}
    \label{fig:cutoff_vs_usage}
\end{figure}

\subsubsection{Dynamic Computation Allocation}
\label{sec:dynamic_computation_allocation}


\begin{figure*}[t]
    \centering
    \begin{minipage}{0.38\textwidth}
        \centering
        \includegraphics[width=\textwidth]{figures/expert_specialization/passage_D_gsm8k_0.pdf}
        \subcaption{Per-token expert routing on a GSM8K passage.}
    \end{minipage}
    \hfill
    \begin{minipage}{0.58\textwidth}
        \centering
        \includegraphics[width=\textwidth]{figures/expert_specialization/heatmap_2k_warmup.pdf}
        \subcaption{Expert activation heatmap. \textbf{Top:} EC with batch size 2k shows less specialization. \textbf{Bottom:} GEC with warmup shows more extreme activation patterns, suggesting more domain-aware routing.}
    \end{minipage}
    % \caption{Expert specialization analysis. (a)~Token-level routing visualization on GSM8K passage. The GEC+warmup 
    % model chooses to do assign more computation to the beginning of sentences. (b)~Activation frequency across 
    % domains reveals clear expert specialization patterns under GEC routing.}
    \caption{Expert specialization analysis. (a)~Token-level activation intensity on a GSM8K passage, colored by total fanout (sum of experts activated across layers). The model assigns more computation to structurally important tokens (punctuation, sentence boundaries, numerical results) than to common content words. (b)~Expert token ratio heatmaps for HumanEval (code) and GSM8K (math). \textbf{Top:} EC (batch size 2k). \textbf{Bottom:} GEC with warmup. GEC achieves sharper patterns in expert activation, suggesting more domain-aware routing and specialization.}
    \label{fig:expert_specialization}
\end{figure*}

A key advantage of GEC/EC is its ability to dynamically allocate computation to different tokens. We visualize this dynamic allocation at the token level. Figure~\ref{fig:expert_specialization}(a) shows per-token activation intensity on a GSM8K passage~\citep{cobbe2021training}, colored by total fanout (the sum of experts activated across layers). The model allocates more computation to structurally important tokens---such as punctuation, sentence boundaries, and numerical results---while assigning fewer experts to common content words. This adaptive allocation suggests that EC/GEC learns to concentrate capacity on tokens requiring more complex reasoning. Additional visualizations are in Appendix~\ref{app:additional_results} (Figures~\ref{fig:passage_fanout}--\ref{fig:passage_intensity}).

\subsubsection{Expert Specialization}
\label{sec:expert_specialization} 


We follow Global LBL~\cite{qiu2025demonsdetailimplementingload} to evaluate expert specialization across EC with various batch sizes (2k, 8k, 64k, 512k) and GEC. For each configuration, we measure the expert token ratio---the fraction of tokens from a given domain routed to each expert---across HumanEval~\citep{chen2021evaluating} (code) and GSM8K~\citep{cobbe2021training} (math) evaluation sets. Figure~\ref{fig:expert_specialization}(b) compares EC (batch size 2k) with GEC (warmup). Both exhibit clear specialization: certain experts consistently attract domain-specific tokens, visible as concentrated dark cells in the heatmap. GEC achieves specialization comparable to EC without requiring large-batch coordination at inference. The full comparison across all batch sizes (Appendix~\ref{app:additional_results}, Figure~\ref{fig:heatmap_combined}) shows that EC specialization sharpens with larger batches---patterns become more concentrated from 2k to 512k---while GEC warmup matches the large-batch EC pattern.



\subsubsection{Capacity Constraints}
\label{sec:capacity}

GEC enforces capacity constraints during training: each expert processes between $(1-\text{capacity\_factor}) \cdot N / E$ and $(1+\text{capacity\_factor}) \cdot N / E$ tokens per batch ($\text{capacity\_factor} = 0.5$), with excess tokens dropped or capacity padded. Since these constraints are absent at inference, frequent triggering would cause train-inference mismatch. Figure~\ref{fig:capacity_constraints} shows that \textbf{capacity constraints are triggered infrequently}: after warmup, both saturation and starvation rates remain low. This confirms that train-inference mismatch from capacity constraints is minimal.

\subsection{Comparison to Token Choice}
\label{sec:tc_comparison}

\begin{figure}[h]
    \centering
    \includegraphics[width=0.9\linewidth]{figures/figures/analysis/load_balancing_vs_deepseek/L1_cutoff_ema_deepseek_vs_ec.pdf}
    \caption{Layer-1 cutoff-EMA (expert 0) under EC/GEC vs DeepSeek loss-free load balancing.}
    \label{fig:l1_cutoff_ema_deepseek_vs_ec}
\end{figure}

\begin{figure*}[t]
    \centering
    \begin{tabular}{@{}ccc@{}}
        \includegraphics[width=0.32\textwidth]{figures/figures/ablations/warmup/warmup_a_cutoff.pdf} &
        \includegraphics[width=0.32\textwidth]{figures/figures/ablations/warmup/warmup_b_raw_usage.pdf} &
        \includegraphics[width=0.32\textwidth]{figures/figures/ablations/warmup/warmup_c_starvation.pdf} \\
        (a) L9 cutoff vs EMA & (b) Raw expert usage & (c) Starvation rate \\[6pt]
        \includegraphics[width=0.32\textwidth]{figures/figures/ablations/warmup/warmup_d_cutoff_ema.pdf} &
        \includegraphics[width=0.32\textwidth]{figures/figures/ablations/warmup/warmup_e_logits_std.pdf} &
        \includegraphics[width=0.32\textwidth]{figures/figures/ablations/warmup/warmup_f_gate_std.pdf} \\
        (d) L6 cutoff-EMA & (e) Router logits std & (f) Gate std \\
    \end{tabular}
    \caption{Effect of TopK warmup on GEC training dynamics (first 8k steps). Before 4k steps, GEC without warmup exhibits unstable threshold routing: (a) the cutoff-EMA lags behind the actual cutoff, (b) raw expert usage is low, and (c) starvation rate is high as the capacity lower bound is frequently triggered. With warmup, the cutoff-EMA trajectory stabilizes (d), and router outputs show lower variance in both logits (e) and gates (f). Note: \texttt{ec\_shared\_bsz512k} does not log raw usage and underflow metrics, so panels (b) and (c) show only the two GEC runs.}
    \label{fig:warmup}
\end{figure*}

EC/GEC achieves more stable load balancing than Token Choice, especially in early layers. Figure~\ref{fig:l1_cutoff_ema_deepseek_vs_ec} compares cutoff-EMA (expert 0) at layer 1: EC/GEC quickly stabilizes, while DeepSeek's loss-free controller drifts upward, indicating persistent imbalance. This is consistent with reports that loss-free load balancing becomes unstable when load statistics are noisy~\cite{wang2024auxiliarylossfreeloadbalancingstrategy}. We conjecture that once routing concentrates, bias updates lose their corrective signal.

\subsection{Ablations}

\paragraph{Warmup}

We find warmup crucial for GEC. In the early stages of training, the cutoff threshold is not yet stable, while the EMA lags behind the actual cutoff threshold because of slow update speed ($1/(1-\beta) \approx 1000$ steps). As a result, the threshold-based routing becomes unreliable: tokens that should be routed are dropped, and the capacity lower bound is frequently triggered (Figure~\ref{fig:warmup}c). This leads to undertrained experts during early training. To address this, we warm up the routing by using TopK selection for the first 4,000 steps before switching to threshold-based routing. As shown in Figure~\ref{fig:warmup}, warmup stabilizes the cutoff-EMA trajectory (a, d), increases raw expert usage (b), and reduces starvation rate (c). We also observe that without warmup, the router exhibits higher variance in both logits (e) and gate outputs (f), suggesting less stable gradient signals.

\paragraph{Shared Expert.}

We report results of EC and GEC with and without shared expert. In both cases, shared expert improves loss by roughly 0.02. We suspect while later layers needs early layers to impower the router, sometimes early layers have no activated experts, causing ineffective routing. See Table~\ref{tab:shared_expert_ablation} for more details.

\begin{table}[h]
\centering
\caption{Ablation on the shared expert mechanism. In both GEC and EC, shared expert improves loss by roughly 0.02. }
\label{tab:shared_expert_ablation}
\small
\setlength{\tabcolsep}{4pt}
\begin{tabularx}{\columnwidth}{@{}>{\raggedright\arraybackslash}X c r r@{}}
\toprule
\textbf{Method} & \textbf{Shared} & \textbf{CE} & \textbf{CORE} \\
\midrule
EC (bsz 512k) & Yes & 2.843 & 19.94 \\
EC (bsz 512k) & No  & 2.862 & 16.307\\
\midrule
GEC ($\beta{=}0.999$) & Yes & 2.844 & 16.867 \\
GEC ($\beta{=}0.999$) & No  & 2.862 & 18.515 \\
\bottomrule
\end{tabularx}
\end{table}



% \paragraph{EMA Decay.} \textcolor{red}{waiting for a 0.99 warmup run.}

\section{Conclusion}

We introduce Global Expert Choice (GEC), a routing mechanism that resolves the fundamental causality issue in Expert Choice (EC) models while preserving their load-balancing advantages. By maintaining an exponential moving average of each expert's selection threshold—estimated from historical batches rather than within-batch top-$k$ selection—GEC enables fully causal routing: each token's routing decision depends only on past statistics, eliminating the need for future token access at both training and inference time. Our experiments demonstrate that GEC achieves competitive performance with EC (matching validation loss at 2.84) while outperforming Token Choice by 0.067 in cross-entropy loss, all while enabling causal autoregressive generation. The cutoff-EMA mechanism provides stable routing thresholds that accurately approximate EC's top-$k$ boundaries, as evidenced by the minimal train-inference gap observed across all metrics. We further show that a warmup strategy—using EC's top-$k$ routing before transitioning to threshold-based selection—stabilizes early training dynamics. These findings suggest that the perceived incompatibility between Expert Choice routing and causal language modeling can be effectively bridged through population-level threshold estimation, opening new directions for scalable MoE architectures.

% \section*{Acknowledgments}
% Use unnumbered first level headings for the acknowledgments. All
% acknowledgments, including those to funding agencies, go at the end of the paper.

\section*{Impact Statement}

This paper presents work whose goal is to advance the field of Machine Learning. There are many potential societal consequences of our work, none which we feel must be specifically highlighted here.

\bibliography{example_paper}
\bibliographystyle{icml2026}


\appendix

\section{Constraint Optimization View}
\label{app:constraint_view}

To understand the theoretical motivation behind Global Expert Choice, it is helpful to formulate MoE routing as a constrained optimization problem. Let $r \in \mathbb{R}^{N \times GE}$ be the router scores (logits) for a batch of $N$ tokens and $GE$ experts. We seek an assignment matrix $z \in \{0, 1\}^{N \times GE}$ that maximizes the total routing score subject to computational constraints.

\paragraph{The Primal Problem.}
The standard routing goal is to maximize the sum of scores for selected experts, subject to two key constraints:
\begin{equation}
\label{eq:primal}
\begin{aligned}
\max_{z} \quad & \sum_{t=1}^N \sum_{i=1}^{GE} z_{t,i} r_{t,i} \\
\text{s.t.} \quad & \sum_{i=1}^{GE} z_{t,i} = G, \quad \forall t \quad \text{(Sparsity)} \\
& \sum_{t=1}^N z_{t,i} = k, \quad \forall i \quad \text{(Load Balancing)} \\
& z_{t,i} \in \{0, 1\}
\end{aligned}
\end{equation}
Here, the \textbf{Sparsity} constraint ensures each token selects exactly $G$ experts (Token Choice). The \textbf{Load Balancing} constraint ensures each expert processes exactly $k = \frac{N}{E}$ tokens.

\paragraph{Difficulty of Top-$G$ Routing.}
Solving \eqref{eq:primal} exactly requires combinatorial algorithms like the Hungarian algorithm. This is computationally expensive ($O(N^3)$) and non-differentiable.
Standard Token Choice (TC) routing strictly enforces the Sparsity constraint ($z_{t,i} = 1 \iff i \in \text{Top}_G(r_{t, \cdot})$) but violates the Load Balancing constraint. Methods like Auxiliary Loss or Capacity Clipping are merely heuristics to approximate the Load Balancing constraint while keeping the Sparsity constraint strict.

\paragraph{Batch Expert Choice: Relaxing the Sparsity Constraint.}
Expert Choice (EC)~\cite{zhou2022mixture} simplifies the problem by relaxing the strict Sparsity constraint. If we drop the requirement that every token must pick exactly $G$ experts, and instead only enforce the Load Balancing constraint (perfect load balancing), the optimization problem decouples across experts:
\begin{equation}
\label{eq:ec_relax}
\max_{z} \sum_{i=1}^{GE} \left( \sum_{t=1}^N z_{t,i} r_{t,i} \right) \quad \text{s.t.} \quad \sum_{t=1}^N z_{t,i} = k, \forall i
\end{equation}
The optimal solution is trivial: each expert $i$ independently selects the top-$k$ tokens with the highest scores $r_{t,i}$.
\begin{equation}
z_{t,i} = \mathbb{1}\{t \in \text{Top}_k(r_{\cdot, i})\}
\end{equation}
This \textbf{Batch Expert Choice} guarantees perfect load balancing and maximizes utility under the load balancing constraint. However, the decision $z_{t,i}$ now depends on the scores of all other tokens in the batch (to determine the rank of $t$), which introduces non-causal dependencies in autoregressive generation.

\paragraph{Global Expert Choice: Stochastic Relaxation.}
Global Expert Choice (GEC) resolves the causality issue by further relaxing the Load Balancing constraint from a hard, per-batch constraint to a \textbf{stochastic expectation}:
\begin{equation}
\mathbb{E}_{\text{data}} \left[ \sum_{t=1}^N z_{t,i} \right] = k
\end{equation}
Instead of enforcing exactly $k$ tokens per batch, we require that experts receive $k$ tokens \textit{on average} over the data distribution. This allows us to replace the batch-dependent ranking with a population-level thresholding:
\begin{equation}
z_{t,i} = \mathbf{1}\{r_{t,i} > c_i\}
\end{equation}
where the threshold $c_i$ is the $(1 - \frac{1}{E})$-quantile of the expert's score distribution $P(r_{\cdot, i})$. By estimating $c_i$ from historical statistics (via EMA), GEC decouples the routing decision: $z_{t,i}$ depends only on $r_{t,i}$ and the global threshold $c_i$, fully restoring causality while satisfying load balancing in expectation.

\section{Future Information Leakage for Expert Choice Models}
\label{app:future_leakage_ec}

In DeepSeek's loss-free load balancing paper~\citep{wang2024auxiliarylossfreeloadbalancingstrategy}, they give an upper bound on the future information leakage of Expert Choice (EC) to be superlinear in the number of tokens. They considered all potential selection combinations $\binom{N}{k}$ for choose $k=N/E$ tokens out of $N$ tokens, which makes upper bound $\log_2\binom{N}{k} = O(N \log N)$.

We consider two scenarios: when cutoff threshold is expressed as a finite precision float, it is trivial that the total future information leakage is at most the number of bits to represent the cutoff threshold. However, when cutoff threshold is of infinite precision, we show that we can indeed leak at least $O(N \log N)$ bits of future information, making the bound tight.

We arrange the following sessions as follows:

\begin{enumerate}
    \item First, we provide a formal definition of future information leakage.
    \item Then, we show for finite precision, the total leakage is constant, which means per token leakage is 0 as batch size increases.
    \item Then, we show for infinite precision, we describe an encoding strategy that can leak at least $O(N \log N)$ bits of future information, making the bound tight. The idea is that we can break the cutoff space into $2^N$ small intervals, and injectively (though not surjectively) map each potential selection combination to a unique interval in a memoryless way. 
    \item We formally prove that the encoding strategy can leak at least $O(N \log N)$ bits of future information.
\end{enumerate}

Gladly, since we rely on finite precision float to represent the cutoff threshold, GEC is still causal. 

\subsection{Definition of Future Information Leakage}

\begin{definition}[Future information leakage]
Fix a sequence length $N$ and a deterministic selection rule $F$ that maps a logit sequence $r_{1:N}$ to a subset $F(r_{1:N}) \subseteq [N]$ (where $[N] \triangleq \{1,\dots,N\}$). Here $r_{1:N}$ denotes the entire sequence of router scores/logits across tokens, and we write $z_t \triangleq \mathbf{1}\!\left[t\in F(r_{1:N})\right]$ for the induced selection indicator.

An \emph{advice variable} is any function $A=\alpha(r_{1:N})$ with finite range, where $\alpha$ is an \emph{encoder} that maps the full logit sequence to a finite label. We write $\mathrm{Range}(\alpha)\triangleq \{\alpha(r_{1:N}) : r_{1:N}\}$ for the set of labels that can be produced by $\alpha$.

We say $A$ \emph{causalizes} $F$ if there exist functions $\{g_t\}_{t=1}^N$ such that for all $r_{1:N}$ and all $t\in[N]$,
\[
z_t \;=\; g_t\!\left(r_{1:t}, A\right).
\]
The future information leakage of $F$ on length $N$ is
\[
\mathcal{L}_F([1\!:\!N]) \;\triangleq\; \min_{\alpha,\,\{g_t\}} \log_2\!\left|\mathrm{Range}(\alpha)\right|.
\]
\end{definition}

\subsection{Finite-precision cutoff implies constant leakage}

With this definition, it is trivial to show that, under finite precision float representation, the total future information leakage is constant, which means per token leakage is 0 as batch size increases.

In our case, $F$ is the Expert Choice selection rule induced by a cutoff threshold: for each expert, tokens are selected by comparing their router scores against the expert's cutoff (equivalently, selecting those above the cutoff, which matches top-$k$ when the cutoff is set to the $k$-th order statistic).

\begin{theorem}[Finite-precision cutoff implies constant leakage]
If the cutoff threshold is represented with $b$ bits of precision (e.g., $b=16$ for \texttt{bf16} or $b=32$ for \texttt{fp32}), then the future information leakage satisfies $\mathcal{L}_F([1\!:\!N]) \le b$ for all $N$.
\end{theorem}

\begin{proof}
This is an upper bound via a particular advice choice. Let the advice be the cutoff itself: $A=\alpha(r_{1:N}) \triangleq \beta$, encoded in $b$ bits, so $|\mathrm{Range}(\alpha)| \le 2^b$. Given $A=\beta$, the selection indicator at time $t$ is a causal function of the prefix (in fact, of $r_t$) and $\beta$ by thresholding, hence $A$ causalizes $F$. Therefore $\mathcal{L}_F([1\!:\!N]) \le \log_2 |\mathrm{Range}(\alpha)| \le b$.
\end{proof}

\subsection{Upper bound on Infinite-precision future information leakage}

The strategy we used in the previous subsection does not work for infinite precision cutoff, as it takes infinite bits of communication to represent the cutoff. Thus, we need a different encoding strategy.

Previously, Loss-free Load Balancing paper~\citet{wang2024auxiliarylossfreeloadbalancingstrategy} gave a combinatorial upper bound on the information carried by an Expert Choice allocation when we allow all admissible token-to-expert assignments consistent with the sparsity pattern.

Let $N$ denote the number of tokens in the routing pool and $GE$ the number of experts. Each token activates $G$ experts, so the MoE sparsity is $\frac{1}{E}$. For an MoE layer in Expert Choice, the maximum information leakage $L$ (bits per token) is:
\begin{align}
L
&= \frac{GE}{N}\log_2\!\binom{N}{\frac{N}{E}} \nonumber\\
&> \frac{GE}{N}\cdot\frac{N}{E}\log_2(E-1) \nonumber\\
&= G\log_2(E-1).
\tag{7}
\end{align}

For a model with sparsity $\frac{1}{E}=\frac{2}{16}=0.125$ and 9 MoE layers, the total leakage information is more than 50 bits per token.

\subsection{Encoding Strategy and Decoding Procedure}

The above upper bound is pretty intuitive to understand. We ask, can we reach this bound? Surprisingly, we can. In this subsection, we describe an encoding strategy where we create an injective mapping from the set of all possible expert choice combinations to the range of the infinite precision cutoff.

Suppose the cutoff $c \in [0,1]$. We partition this space into $2^N$ dyadic intervals. Let $z \in \{0,1\}^N$ be the target expert choice combination (codeword). We map $z$ to the interval:
\[
I(z)\;\triangleq\;\Big[\sum_{t=1}^N (1-z_t)2^{-t},\ \sum_{t=1}^N (1-z_t)2^{-t} + 2^{-N}\Big).
\]
This orders codewords in descending order: $00\cdots 00$ maps to the rightmost interval and $11\cdots 11$ maps to the leftmost. Choose a cutoff value $c\in I(z)$. Obviously, this mapping is injective, i.e. each codeword maps to a unique interval. However, since not all combinations are possible, it is not surjective.

\paragraph{Information from the Past.}
For a token $t$, what information about the past routing logits $r_{1:t}$ and decisions $z_{1:t}$ can we use to guide the selection of the interval? Turns out, the only information we get is the current bracket $[\ell, u)$ containing the decision boundary. The upper bound $u$ is determined by the lowest logit of the selected tokens, and the lower bound $\ell$ by the highest logit of the unselected tokens. To maximize information leakage, we minimize help from the past by always choosing the next query logit to be in the middle of the interval, $r_t = (\ell+u)/2$. Then, the routing decision $z_t = \mathbf{1}\{r_t \ge c\}$ reveals exactly one bit of the cutoff, requiring full future information.

\paragraph{Decoding Procedure.}
We formalize this procedure in Algorithm~\ref{alg:cutoff_decode}.

\begin{algorithm}[t]
\caption{Binary-search decoding of an infinite-precision cutoff}
\label{alg:cutoff_decode}
\begin{algorithmic}[1]
\STATE \textbf{Input:} horizon $N$; unknown cutoff $c\in[0,1]$; oracle bit $z_t=\mathbf{1}\{r_t\ge c\}$
\STATE $\ell\gets 0,\ u\gets 1$
\FOR{$t=1,\ldots,N$}
\STATE $r_t\gets (\ell+u)/2$ \COMMENT{Query midpoint}
\STATE Observe selection $z_t \in \{0,1\}$
\IF{$z_t=1$}
\STATE $u\gets r_t$ \COMMENT{Selected ($r_t \ge c$) $\implies c \in [\ell, r_t)$}
\ELSE
\STATE $\ell\gets r_t$ \COMMENT{Not selected ($r_t < c$) $\implies c \in [r_t, u)$}
\ENDIF
\ENDFOR
\STATE \textbf{Return} $(z_1,\ldots,z_N)$ and interval $[\ell,u)$
\end{algorithmic}
\end{algorithm}

\begin{figure}[t]
\centering
\includegraphics[width=0.9\linewidth]{figures/illustration/encoding.png}
\caption{Illustration of the binary-search encoding strategy. The cutoff value $c$ partitions the interval $[0,1]$ into regions corresponding to different expert selection patterns, enabling $N$ bits of future information to be encoded in a single real-valued threshold.}
\label{fig:encoding}
\end{figure}

\subsection{Formal Proof}

We now formally prove that the amount of information required to implement the Expert Choice selection causally is lower-bounded by the combinatorial entropy of the selection space. This confirms that the upper bound in Eq.~(7) is tight and that the infinite-precision construction in Algorithm~\ref{alg:cutoff_decode} is optimal in terms of information leakage.

\begin{theorem}
\label{thm:lower_bound}
For a single expert with capacity $k = \lfloor N/E \rfloor$, any causal routing mechanism that can realize all possible top-$k$ assignments requires at least $\log_2 \binom{N}{k}$ bits of non-causal information (advice).
\end{theorem}

\begin{proof}
Let $\mathcal{Z}_k = \{z \in \{0,1\}^N : \sum_{t=1}^N z_t = k\}$ be the set of all valid selection indicators for the expert, with $|\mathcal{Z}_k| = \binom{N}{k}$. We show that distinct advice is necessary for every distinct pattern in $\mathcal{Z}_k$.

Consider the specific family of logit sequences generated by the binary search process in Algorithm~\ref{alg:cutoff_decode}. In this construction, the router logit $r_t$ is the midpoint of the current valid interval $[\ell, u)$, which is determined solely by the history of decisions $z_1, \dots, z_{t-1}$. Consequently, if two selection patterns $z$ and $z'$ share the same prefix $z_{1:t-1}$, they will generate the exact same logit $r_t$ at step $t$.

Suppose there exists an advice encoding with range size strictly less than $\binom{N}{k}$. By the Pigeonhole Principle, at least two distinct valid patterns $z, z' \in \mathcal{Z}_k$ must share the same advice value. Let $t$ be the first index where they differ (i.e., $z_{1:t-1} = z'_{1:t-1}$ but $z_t \neq z'_t$).
Since the prefixes are identical, the generated logit sequence $r_{1:t}$ is identical for both patterns. A causal decoder, which must output a decision at time $t$ based only on $r_{1:t}$ and the advice, receives identical inputs for both cases. It must therefore produce the same output. However, $z_t \neq z'_t$, so the decoder necessarily fails for at least one of the patterns.

Thus, every valid top-$k$ pattern requires a unique advice value. The minimum information leakage is $\log_2 \binom{N}{k}$ bits per expert.
\end{proof}

Summing this lower bound over $GE$ experts recovers the combinatorial quantity in Eq.~(7), proving that Expert Choice routing fundamentally requires significant future information to implement.

\section{Architecture Details}
\label{app:arch_diff}

Our model architecture follows nanochat~\cite{karpathy2025nanochat}, which differs from standard GPT-2~\cite{radford2019language} in several ways. Table~\ref{tab:arch_config} summarizes the key differences.

\begin{table*}[t]
\centering
\caption{Model architecture and size configurations. Architecture features are shared between d12 and d20 (nanochat-style). For MoE variants with $G{=}2$, $E{=}8$: 16 routed experts + 1 shared = 17 total experts. Total params include all expert parameters; active params include only the shared expert plus on average one routed expert per token.}
\label{tab:arch_config}
\small
\begin{tabular}{@{}l|>{\centering\arraybackslash}p{3.5cm}|>{\centering\arraybackslash}p{3cm}|>{\centering\arraybackslash}p{3cm}@{}}
\toprule
\textbf{Feature} & \textbf{GPT-2} & \textbf{d12} & \textbf{d20} \\
\midrule
Tokenization & GPT-2 & \multicolumn{2}{c|}{RustBPE} \\
Vocab Size & 50,257 & \multicolumn{2}{c|}{65,536} \\

Activation & GELU & \multicolumn{2}{c|}{ReLU$^2$~\cite{so2021primer}} \\
Normalization & LayerNorm & \multicolumn{2}{c|}{RMSNorm~\cite{zhang2019rmsnorm}} \\
FFN Dimension & $4 \times d_{\text{model}}$ & \multicolumn{2}{c|}{$4 \times d_{\text{model}}$} \\
Linear Layer Bias & Yes & \multicolumn{2}{c|}{No} \\
Position Encoding & Learned & \multicolumn{2}{c|}{RoPE~\cite{su2024rope}} \\
Head Dimension & 64 & \multicolumn{2}{c|}{128} \\
QK Normalization & No & \multicolumn{2}{c|}{Yes~\cite{dehghani2023vit22b}} \\
Logits Softcapping & No & \multicolumn{2}{c|}{15.0~\cite{team2024gemma2}} \\
Embedding Weights & Tied (wte = lm\_head) & \multicolumn{2}{c|}{Untied} \\

First Layer Dense & --- & \multicolumn{2}{c|}{Yes} \\
\midrule
$n_{\text{embd}}$ & 768 & 768 & 1280 \\
$n_{\text{layer}}$ & 12 & 12 & 20 \\
$n_{\text{head}}$ & 12 & 6 & 10 \\
KV heads & 12 & 2 & 2 \\
Dense Params & 124M & 195M & 561M \\
MoE Total Params & --- & 575M & 2429M \\
MoE Active Params & --- & 195M & 561M \\
\bottomrule
\end{tabular}
\end{table*}

\subsection{Data and Tokenization}

We train on the FineWeb-Edu 100B shuffle dataset~\cite{penedo2024the}, a high-quality educational web corpus. Tokenization uses RustBPE with a vocabulary of 65,536 tokens (64k, power-of-2 aligned for GPU efficiency). We use a sequence length of 2048 tokens.

\subsection{Model Size Configurations}
\label{app:model_sizes}

Table~\ref{tab:arch_config} shows the model configurations used in our experiments. We follow the nanochat naming convention where \texttt{d*} indicates the number of layers, and $n_{\text{embd}} = d \times 64$, with head dimension fixed at 128.

\paragraph{Attention configuration.} We use grouped-query attention~\cite{ainslie2023gqa} with head dimension 128 (larger than GPT-2's 64). The number of attention heads is $n_{\text{head}} = n_{\text{embd}} / 128$, giving 6 heads for d12 and 10 heads for d20. QK normalization is applied before the attention computation.

\paragraph{MoE configuration.} For MoE variants, we use $G{=}2$ (granularity) and $E{=}8$ (expansion), yielding 16 routed experts plus 1 shared expert (17 total). Each expert has dimension $d_{\text{expert}} = 2 \times n_{\text{embd}}$ (half the dense FFN dimension). The shared expert processes every token, while each token activates on average 1 routed expert, matching the dense model's active parameter count.



\section{Training Setup Details}
\label{app:training_setup}



\subsection{Training Hyperparameters}

\begin{table}[h]
\centering
\caption{Training hyperparameters.}
\label{tab:training_hyperparams}
\small
\begin{tabular}{@{}ll@{}}
\toprule
\textbf{Hyperparameter} & \textbf{Value} \\
\midrule
Total tokens & 10B / 11.2B \\
Batch size (tokens) & 524,288 (0.5M) \\
Sequence length & 2048 \\
Muon Warmup steps & No warmup\\
AdamW Warmup steps & 250 \\
Learning rate schedule & Linear decay \\
Min learning rate & $0.1 \times$ peak LR \\
Gradient clipping & None \\
Weight decay & 0.0 \\
AdamW $\beta_1, \beta_2$ & 0.9, 0.95 \\
\bottomrule
\end{tabular}
\end{table}

\subsubsection{Weight Initialization}

We follow the nanochat initialization scheme~\cite{karpathy2025nanochat}, which uses aspect-ratio scaled initialization. For a weight matrix $W \in \mathbb{R}^{d_{\text{out}} \times d_{\text{in}}}$:

\begin{equation}
\text{std} = \frac{1}{\sqrt{d_{\text{in}}}} \cdot \min\left(1, \sqrt{\frac{d_{\text{out}}}{d_{\text{in}}}}\right)
\end{equation}

This formula reduces to standard $1/\sqrt{d_{\text{in}}}$ initialization for square or tall matrices, but scales down variance for wide matrices where $d_{\text{out}} \gg d_{\text{in}}$.

\paragraph{Component-specific initialization:}
\begin{itemize}
    \item \textbf{Embeddings}: $\mathcal{N}(0, 1)$ -- standard normal initialization
    \item \textbf{Output projections} (\texttt{lm\_head}, \texttt{c\_proj}): Zero initialization, critical for Muon optimizer stability
    \item \textbf{Router weights}: $\mathcal{N}(0, 1/\sqrt{d_{\text{in}}})$ -- small init for symmetry breaking
    \item \textbf{Expert weights}: Aspect-ratio scaled for up projections, zero for down projections
    \item \textbf{Attention weights}: Aspect-ratio scaled as above
\end{itemize}

\begin{table}[h]
\centering
\caption{Parameter initialization and optimizer configuration. Aspect-ratio scaled init.\ uses $\text{std} = d_{\text{in}}^{-1/2} \cdot \min(1, \sqrt{d_{\text{out}}/d_{\text{in}}})$. LRs include $\mu$P~\cite{yang2022mup} scaling $\lambda = (d_{\text{model}}/768)^{-1/2}$.}
\label{tab:param_init_lr}
\footnotesize
\begin{tabular}{@{}lccc@{}}
\toprule
\textbf{Parameter} & \textbf{Init.} & \textbf{LR} & \textbf{Opt.} \\
\midrule
$W_E$ (embed.) & $\mathcal{N}(0, 1)$ & $0.2\lambda$ & AdamW \\
$W_{\text{lm}}$ (head) & $\mathbf{0}$ & $0.004\lambda$ & AdamW \\
\midrule
$W_{QKV}$ (attn) & Asp.-ratio & $0.02\lambda$ & Muon \\
$W_O$ (attn proj) & $\mathbf{0}$ & $0.02\lambda$ & Muon \\
\midrule
$W_{\text{router}}$ & $\mathcal{N}(0, d^{-\frac{1}{2}})$ & $0.02\lambda$ & Muon \\
$W^{(e)}_{\uparrow}$ (exp.\ up) & Asp.-ratio & $0.02\lambda$ & Muon \\
$W^{(e)}_{\downarrow}$ (exp.\ dn) & $\mathbf{0}$ & $0.02\lambda$ & Muon \\
$W^{(s)}_{\uparrow}$ (shd.\ up) & Asp.-ratio & $0.02\lambda$ & Muon \\
$W^{(s)}_{\downarrow}$ (shd.\ dn) & $\mathbf{0}$ & $0.02\lambda$ & Muon \\
\bottomrule
\end{tabular}
\end{table}

\subsubsection{Optimizer Configuration}

We use a hybrid optimizer setup following nanochat~\cite{jordan2024muon, liu2025moonlight}:
\begin{itemize}
    \item \textbf{Muon}~\cite{jordan2024muon} for 2D/3D weight matrices (attention, MLP, experts): momentum-based optimizer with Newton-Schulz orthogonalization
    \item \textbf{AdamW}~\cite{loshchilov2019adamw} for embeddings and output head: with learning rate scaling $\propto 1/\sqrt{d_{\text{model}}/768}$
\end{itemize}

No weight decay is used, as Muon provides implicit regularization and language models benefit from memorization.

\subsection{Hardware Infrastructure Details}

\paragraph{Hardware.} We train our models on a single node with 8x NVIDIA B200 GPUs, each with 180GB of memory. 

\paragraph{Code.} For TC models, we use the ScatterMoE backend \citep{tan2024scattermoe}. For EC and GEC models, we write our own custom Pytorch MoE implementation. We use padding to handle variable number of tokens per expert. 

\paragraph{Parallelism.} We rely on Nanochat's implementation of distributed AdamW and Muon optimizer use a ZeRO-2 style gradient synchronization \citep{rajbhandari2020zero}. For EC and GEC models, we write our own expert parallelization all-to-all communication framework. This allows us to use maximum batch size during routing instead of micro-batches, reaching a better usage/cutoff variance trade-off while saving memory.


\section{CORE Evaluation Details}
\label{app:coreeval}

We evaluate using the CORE benchmark~\cite{li2024dclm}, which provides a standardized suite of in-context learning tasks for language model evaluation.

\paragraph{Task types.} The CORE benchmark includes multiple-choice tasks, schema matching tasks, and language modeling tasks, testing various aspects of language understanding.

\paragraph{Metric.} The primary metric is \emph{centered accuracy}, which adjusts for random baseline performance:
\begin{equation}
\text{acc}_{\text{centered}} = \frac{\text{acc} - 0.01 \times \text{baseline}_{\text{random}}}{1.0 - 0.01 \times \text{baseline}_{\text{random}}}
\end{equation}
This normalization ensures that random guessing yields a score near zero, while perfect accuracy yields 1.0. The final CORE Eval score is the mean of centered accuracy across all tasks.

\paragraph{Evaluation protocol.} We evaluate at fixed intervals during training (every 250 steps by default) to track learning dynamics.


\section{Additional Experiment Results}
\label{app:additional_results}

\paragraph{Capacity constraints.}
Figure~\ref{fig:capacity_constraints} shows capacity constraint metrics for GEC (with warmup) from step 4k onward. After warmup, raw expert usage stabilizes around 6.5\%, and both saturation and starvation rates remain low, confirming that capacity constraints are rarely triggered and train-inference mismatch is minimal.

\begin{figure*}[t]
    \centering
    \includegraphics[width=0.95\textwidth]{figures/figures/analysis/ec_vs_gec/capacity_constraints/capacity_metrics.pdf}
    \caption{Capacity constraint behavior during GEC training (from step 4k onward, after warmup). (a) Raw expert usage before capacity capping. (b) Saturation rate: fraction of selected tokens dropped due to capacity limits. (c) Starvation rate: fraction of unused expert capacity. Both saturation and starvation rates remain low, confirming minimal train-inference mismatch.}
    \label{fig:capacity_constraints}
\end{figure*}

This appendix also provides extended expert specialization analysis, complementing the summary in Section~\ref{sec:expert_specialization}.

\paragraph{Per-token routing visualizations.}
Figures~\ref{fig:passage_fanout},~\ref{fig:passage_routing}, and~\ref{fig:passage_intensity} show token-level expert routing for additional passages from GSM8K and HumanEval. Across GSM8K passages, content-bearing tokens---particularly numbers (e.g., ``48'', ``72''), mathematical operators (``/'', ``+'', ``=''), and computation markers (``$<<$'')---consistently receive the highest fanout. Function words and punctuation receive minimal activation, indicating that experts preferentially process semantically rich tokens. In HumanEval passages, a similar pattern holds: code-specific tokens (variable names, operators, keywords) receive higher activation than boilerplate text and whitespace.

\begin{figure*}[t]
    \centering
    \begin{minipage}{0.48\textwidth}
        \centering
        \includegraphics[width=\textwidth]{figures/expert_specialization/passage_A_gsm8k_0.pdf}
        \subcaption{GSM8K.}
    \end{minipage}
    \hfill
    \begin{minipage}{0.48\textwidth}
        \centering
        \includegraphics[width=\textwidth]{figures/expert_specialization/passage_A_humaneval_0.pdf}
        \subcaption{HumanEval.}
    \end{minipage}
    \caption{Per-layer expert fanout on GSM8K and HumanEval passages. Each cell shows the number of experts activated for a given token at a given layer. Numerical and code-specific tokens receive substantially higher fanout than function words.}
    \label{fig:passage_fanout}
\end{figure*}

\begin{figure*}[t]
    \centering
    \begin{minipage}{0.48\textwidth}
        \centering
        \includegraphics[width=\textwidth]{figures/expert_specialization/passage_C_gsm8k_0.pdf}
        \subcaption{GSM8K.}
    \end{minipage}
    \hfill
    \begin{minipage}{0.48\textwidth}
        \centering
        \includegraphics[width=\textwidth]{figures/expert_specialization/passage_C_humaneval_0.pdf}
        \subcaption{HumanEval.}
    \end{minipage}
    \caption{Full expert routing on GSM8K and HumanEval passages. Each panel shows binary expert activation (black = activated) across all layers and experts for every token. Routing patterns reveal domain-specific structure.}
    \label{fig:passage_routing}
\end{figure*}

\begin{figure*}[p]
    \centering
    \includegraphics[width=\textwidth]{figures/expert_specialization/passage_D_humaneval_0.pdf}
    \caption{Token activation intensity on a HumanEval passage. Each token is colored by total fanout (sum of experts activated across all layers). Code-specific tokens (variable names, operators, keywords) receive higher activation than boilerplate text.}
    \label{fig:passage_intensity}
\end{figure*}

\paragraph{Expert activation heatmaps.}
Figure~\ref{fig:heatmap_combined} shows expert token ratios across all routing configurations. Each heatmap plots expert ID (columns) versus layer (rows), with color intensity indicating the fraction of domain-specific tokens routed to each expert. The left column shows HumanEval (code) and the right column shows GSM8K (math).

Several patterns emerge across batch sizes. EC with batch size 2k (top row) shows diffuse activation: while some experts exhibit domain preferences (e.g., concentrated dark cells at specific layer-expert pairs), the overall pattern is noisy with activation spread across many experts. As batch size increases to 8k and 64k, specialization sharpens---dark cells become more concentrated and background activation fades, indicating that experts more consistently capture domain-specific tokens when routing decisions are made over larger token pools. EC at 512k shows the most pronounced specialization, with a small number of experts per layer handling the majority of domain tokens.

GEC with warmup (bottom row) achieves specialization comparable to large-batch EC. The activation patterns closely resemble EC at 512k, with concentrated expert-domain associations across layers. This confirms that GEC's population-level threshold mechanism captures the same routing structure as large-batch top-$k$ selection, without requiring batch size coordination at inference.

Comparing across domains, the HumanEval and GSM8K columns reveal that experts develop \emph{different} specialization patterns for code versus math. Certain experts that are heavily activated for code tokens (e.g., dark cells in the HumanEval column) show low activation for math, and vice versa. This cross-domain differentiation is consistent across all routing configurations, suggesting that expert specialization reflects genuine domain-level structure rather than artifacts of a particular routing strategy.

\begin{figure*}[t]
    \centering
    \includegraphics[width=0.7\textwidth]{figures/expert_specialization/heatmap_2k_to_warmup.pdf}
    \caption{Expert activation heatmaps across routing configurations. Each row corresponds to a routing variant (EC with batch sizes 2k, 8k, 64k, 512k, and GEC with warmup). Columns show HumanEval (code) and GSM8K (math) domains. Color intensity indicates expert token ratio. Specialization sharpens with larger EC batch sizes; GEC warmup achieves comparable patterns without batch size dependence.}
    \label{fig:heatmap_combined}
\end{figure*}



\section{Failed Attempts}

\subsection{Normalization}

Initially, we assumed that dynamic expert count would bring instability in training because of the scale expansion. However, we found that normalization is proven to be ineffective. No norm outperformed fanout norm by 0.04 in CE loss. We suspect that the norm made experts' contribution unpredictable (see Figure~\ref{fig:norm_comparison}).

\begin{figure}[h]
    \centering
    \includegraphics[width=0.8\linewidth]{figures/figures/failed_experiments/norm/eval_loss_norm_comparison.pdf}
    \caption{Comparison of evaluation loss with and without normalization. The configuration without normalization (blue) consistently achieves lower loss than the fanout-normalized variant (orange).}
    \label{fig:norm_comparison}
\end{figure} 











\end{document}


% This document was modified from the file originally made available by
% Pat Langley and Andrea Danyluk for ICML-2K. This version was created
% by Iain Murray in 2018, and modified by Alexandre Bouchard in
% 2019 and 2021 and by Csaba Szepesvari, Gang Niu and Sivan Sabato in 2022.
% Modified again in 2023 and 2024 by Sivan Sabato and Jonathan Scarlett.
% Previous contributors include Dan Roy, Lise Getoor and Tobias
% Scheffer, which was slightly modified from the 2010 version by
% Thorsten Joachims & Johannes Fuernkranz, slightly modified from the
% 2009 version by Kiri Wagstaff and Sam Roweis's 2008 version, which is
% slightly modified from Prasad Tadepalli's 2007 version which is a
% lightly changed version of the previous year's version by Andrew
% Moore, which was in turn edited from those of Kristian Kersting and
% Codrina Lauth. Alex Smola contributed to the algorithmic style files.
